\documentclass[11pt]{article}

\usepackage[a4paper,margin=1in]{geometry}
\usepackage{amsmath,amssymb}
\usepackage{hyperref}

\title{Mathematical Model for the Exhaust-Plume Study Solver}
\author{Exhaust-Plume project}
\date{\today}

\newcommand{\Mach}{\mathit{M}}
\newcommand{\gammaa}{\gamma}
\newcommand{\Rgas}{R}
\newcommand{\dd}{\,\mathrm{d}}

\begin{document}
\maketitle

\begin{abstract}
This note records the equations and modeling assumptions used by the
standalone compressible exhaust-plume study solver.  The implementation is a
calorically perfect, ideal-gas model that combines one-dimensional nozzle
relations with Prandtl--Meyer expansion fans, oblique shocks, and a
two-dimensional geometric construction of plume zones.  It is intended for
reproducible exploratory studies rather than certification analysis.
\end{abstract}

\section{Assumptions and notation}

The gas is modeled as calorically perfect with constant ratio of specific heats
\(\gamma=C_p/C_v\).  The equation of state and speed of sound are
\begin{equation}
  p = \rho \Rgas T,
  \qquad
  a = \sqrt{\gammaa \frac{p}{\rho}},
  \qquad
  V = \Mach a .
\end{equation}
Here \(p\) is static pressure, \(\rho\) is density, \(T\) is static
temperature, \(a\) is speed of sound, \(V\) is flow speed, and \(\Mach\) is
Mach number.  The code uses SI units: pascals, kelvin, kilograms per cubic
meter, meters, and seconds.  Angles are stored in degrees in the public flow
state classes and converted to radians for trigonometric operations.

\section{Isentropic relations}

For an isentropic change between a static state and its stagnation state,
\begin{align}
  \frac{T_0}{T} &= 1 + \frac{\gammaa-1}{2}\Mach^2, \\
  \frac{p_0}{p} &=
    \left(1 + \frac{\gammaa-1}{2}\Mach^2\right)^{\gammaa/(\gammaa-1)}, \\
  \frac{\rho_0}{\rho} &=
    \left(1 + \frac{\gammaa-1}{2}\Mach^2\right)^{1/(\gammaa-1)}.
\end{align}
The nozzle-exit helper takes \(p_0\), \(T_0\), \(\Mach\), and \(\gammaa\),
then evaluates these relations to obtain the exit static state.  Its current
gas constant is based on the configured dry-air molar mass.

For a quasi-one-dimensional nozzle, the area--Mach relation is
\begin{equation}
  \frac{A}{A^*} = \frac{1}{\Mach}
  \left[\frac{2}{\gammaa+1}
  \left(1+\frac{\gammaa-1}{2}\Mach^2\right)\right]^{{(\gammaa+1)}/{(2(\gammaa-1))}},
\end{equation}
where \(A^*\) is the sonic throat area.  The study solver uses the
supersonic branch for an exit Mach number greater than one.

\section{Prandtl--Meyer expansion}

For a supersonic expansion, the Prandtl--Meyer function is
\begin{equation}
  \nu(\Mach) =
  \sqrt{\frac{\gammaa+1}{\gammaa-1}}
  \tan^{-1}\!\left(\sqrt{\frac{\gammaa-1}{\gammaa+1}(\Mach^2-1)}\right)
  - \tan^{-1}\!\left(\sqrt{\Mach^2-1}\right).
\end{equation}
The flow turning angle between two states is
\begin{equation}
  \Delta\theta = \nu(\Mach_2)-\nu(\Mach_1).
\end{equation}
The expansion fan preserves stagnation pressure and temperature while changing
the static state according to the isentropic relations.

\section{Normal and oblique shocks}

An oblique shock is reduced to a normal shock using the upstream Mach number
normal to the shock,
\begin{equation}
  \Mach_{n1} = \Mach_1\sin\beta,
\end{equation}
where \(\beta\) is the shock angle.  The normal-shock relations used by the
solver are
\begin{align}
  \frac{p_2}{p_1} &= 1 + \frac{2\gammaa}{\gammaa+1}(\Mach_{n1}^2-1), \\
  \frac{\rho_2}{\rho_1} &=
    \frac{(\gammaa+1)\Mach_{n1}^2}
         {2+(\gammaa-1)\Mach_{n1}^2}, \\
  \frac{T_2}{T_1} &= \frac{p_2/p_1}{\rho_2/\rho_1}, \\
  \Mach_{n2}^2 &=
    \frac{1+\frac{\gammaa-1}{2}\Mach_{n1}^2}
         {\gammaa\Mach_{n1}^2-\frac{\gammaa-1}{2}}.
\end{align}
The downstream Mach number for the oblique shock is reconstructed as
\begin{equation}
  \Mach_2 = \frac{\Mach_{n2}}{\sin(\beta-\theta)},
\end{equation}
where \(\theta\) is the flow-turning angle.

The shock angle is determined from the \(\theta\)--\(\beta\)--\(\Mach\)
relation,
\begin{equation}
  \tan\theta = 2\cot\beta\,
  \frac{\Mach_1^2\sin^2\beta-1}
       {\Mach_1^2\left(\gammaa+\cos(2\beta)\right)+2}.
\end{equation}
The weak-shock branch is used for ordinary compression-wave construction.

\section{Plume-zone construction}

The solver first constructs the nozzle-exit zone from total conditions.  Let
\(p_e\) be its static pressure and \(p_a\) the atmospheric pressure.

\begin{itemize}
  \item If \(p_e > p_a\), the plume is underexpanded.  The solver constructs
    an expansion-fan sequence, reflects the characteristics across the
    centerline, and adds oblique compression waves to approximately equalize
    the plume pressure with the atmosphere.
  \item If \(p_e \leq p_a\), the plume is treated as overexpanded.  An initial
    pair of oblique-shock precursor zones raises the pressure before the
    reflected expansion/compression construction is applied.
\end{itemize}

The geometric construction uses straight characteristic segments, line
intersections, and a fitted parabola for one reflected slip-line boundary.
Some intermediate regions are intentionally represented as approximate study
zones; the method is not a full viscous or reacting-flow solution.

\section{Conservation checks}

The benchmark tests in the repository verify the algebraic relations above.
For a normal shock they additionally check conservation of mass flux,
streamwise momentum flux, and total enthalpy:
\begin{align}
  \rho_1 V_1 &= \rho_2 V_2, \\
  p_1+\rho_1 V_1^2 &= p_2+\rho_2 V_2^2, \\
  C_pT_1+\frac{V_1^2}{2} &= C_pT_2+\frac{V_2^2}{2}.
\end{align}
For isentropic transformations, the tests verify that converting static
properties to stagnation properties and back preserves the supplied state.

\section{Known limits}

The reflected geometry can become ill-conditioned for extreme pressure ratios
and very low or very high Mach numbers.  The current implementation reports
invalid input values early, but it does not claim convergence for every
physically imaginable input.  Future work should replace the geometric hacks
with explicit characteristic validity checks and report a structured solver
status alongside the zone list.

\end{document}
